\section{Related Work}
\label{sec:related}

\paragraph{Financial NLP.}
Loughran and McDonald~\shortcite{loughran2011liability} flagged
systematic miscalibration of generic sentiment lexicons in financial
text; Financial PhraseBank~\cite{malo2014fpb} grew out of that line.
FinBERT~\cite{araci2019finbert,yang2020finbert} is the de facto open
domain transformer; surveys~\cite{wang2024finllmsurvey} document a
recent wave of billion-parameter financial reasoners
(BloombergGPT~\cite{wu2023bloomberggpt},
FinGPT~\cite{yang2023fingpt}, PIXIU~\cite{xie2023pixiu}). We ask
\emph{not} which model is best, but how small each agent in a
committee can become.

\paragraph{Multi-agent LLMs.}
Query-by-committee~\cite{settles2009al} and ensemble
disagreement~\cite{lakshminarayanan2017ensembles} are the classical
ancestors of our SDI. Recent multi-agent LLM frameworks ---
AutoGen~\cite{wu2023autogen}, MetaGPT~\cite{hong2023metagpt},
CAMEL~\cite{li2023camel}, multi-agent
debate~\cite{du2023debate,liang2023encouraging}, LLM-as-a-judge
\cite{zheng2023judging}, self-consistency
\cite{wang2023selfconsistency} --- typically use the \emph{same}
large model for every agent. Our negative result on same-size
persona vote (Section~\ref{sec:experiments:scaling}) shows this does
not substitute for the granularity-stratified diversity that drives
our critic plateau.

\paragraph{Cost-aware routing and semantic caching.}
FrugalGPT~\cite{chen2023frugalgpt}, RouteLLM~\cite{ong2024routellm},
and mixture-of-models routing~\cite{wang2022mixture,shnitzer2023llmrouting}
cascade between cheap and expensive models per query. Semantic
caches~\cite{wang2024caching} (a sibling of our SCD) reduce inference
cost by reusing answers for near-duplicate queries; our SCD differs
in that it caches \emph{committee decisions} (not single-model
outputs) and uses multilingual sentence
embeddings~\cite{reimers2019sbert,reimers2020multilingualkd} for
cross-lingual canonicalisation.
